% Part: counterfactuals
% Chapter: introduction
% Section: material-conditional

\documentclass[../../../include/open-logic-section]{subfiles}

\begin{document}

\olfileid{cnt}{int}{mat}

\olsection{Kondisional Material}

Ing wujud kang paling prasaja ing basa Inggris, kondisional iku
ukara awujud ``If \dots then \dots'' (``Yen ... mula ...''), kanthi
saben \dots{} minangka ukara dhewe. Tuladhane, yen diterjemahake,
``Yen pelayan omah iku kang nindakake, mula tukang kebone ora
luput.'' Ing kuliah pambuka logika, kita sinau nyimbolake
kondisional nganggo konektif $\lif$: perangan kang dituduhake
dening \dots{} disimbolake, upamane nganggo !!{formula}s $!A$
lan~$!B$, banjur kabeh kondisional disimbolake kanthi $!A \lif !B$.

Konektif $\lif$ iku \emph{fungsional tumrap bebener}, yaiku nilai
bebener---$\True$ utawa $\False$---saka $!A \lif !B$ ditemtokake
dening nilai bebener saka $!A$ lan~$!B$: $!A \lif !B$ bener yen
lan mung yen $!A$~luput utawa $!B$~bener, lan luput ing kasus
liyane. Tumrap penetapan nilai bebener~$\pAssign v$, kita nemtokake
$\pSat{v}{!A \lif !B}$ yen lan mung yen $\pSat/{v}{!A}$ utawa
$\pSat{v}{!B}$. Konektif $\lif$ kanthi semantik iki diarani
\emph{kondisional material}.

Definisi iki ngasilake sawetara fakta logis dhasar. Kapisan,
teorema deduksi lumaku kanggo kondisional material:
\begin{align}
  \text{Yen } \Gamma, !A \Entails !B \text{ mula } \Gamma \Entails !A \lif !B
\end{align}
Kondisional iku fungsional tumrap bebener: $!A \lif !B$ lan
$\lnot !A \lor !B$ ekuivalen:
\begin{align}
  !A \lif !B & \Entails \lnot !A \lor !B\\
  \lnot !A \lor !B & \Entails !A \lif !B
  \intertext{Kondisional material minangka konsekuensi semantis saka
    konsekuene lan saka negasi antesedene:}
  !B & \Entails !A \lif !B\\
  \lnot !A & \Entails !A \lif !B
  \intertext{Kondisional material kang luput ekuivalen karo konjungsi
    antesedene lan negasi konsekuene: yen $!A \lif !B$ luput,
    $!A \land \lnot !B$ bener, lan kosok balene:}
  \lnot(!A \lif !B) & \Entails !A \land \lnot !B\\
  !A \land \lnot !B & \Entails \lnot(!A \lif !B)
  \intertext{Kondisional material ndhukung modus ponens:}
  !A, !A \lif !B & \Entails !B
  \intertext{Konsekuen-konsekuen saka kondisional material kanthi
    anteseden kang padha bisa digandhengake:}
  !A \lif !B,  !A \lif !C & \Entails !A \lif (!B \land !C)
  \intertext{Kita tansah bisa nguwatake anteseden, yaiku kondisional
    iku \emph{monoton}:}
  !A \lif !B & \Entails (!A \land !C) \lif !B
  \intertext{Kondisional material transitif, yaiku aturan rantai sah:}
  !A \lif !B, !B \lif !C & \Entails !A \lif !C
  \intertext{Kondisional material ekuivalen karo kontrapositife:}
  !A \lif !B & \Entails \lnot !B \lif \lnot !A\\
  \lnot !B \lif \lnot !A & \Entails !A \lif !B
\end{align}

Kabeh iki inferensi kang migunani lan ora nuwuhake masalah ing
penalaran matematika. Nanging pustaka filsafat lan linguistik kebak
tuladha kang diajokake minangka panyanggah tumrap inferensi kang
padha ing konteks dudu matematika. Tuladha-tuladha iku nuduhake yen
kondisional material~$\lif$ dudu---utawa paling ora ora tansah
dadi---konektif kang trep kanggo nyimbolake pratelan basa Inggris
``if \dots then \dots''.

\end{document}
